% Emacs settings: -*-mode: latex; TeX-master: "manual.tex"; -*-

\chapter{Instrument examples}
\label{s:instrument}

In this section, we present a few typical instruments from the \MCS example
library.  We then give a longer description of three selected worked
examples: a simple sample test instrument (\ref{s:V-instr}), the historical
Ris\o\ triple-axis spectrometer TAS1 (\ref{s:TAS1}), and the ISIS
time-of-flight spectrometer PRISMA (\ref{s:PRISMA}).

These instrument files are included in the \MCS distribution in the
\verb+examples/+ directory (\verb+mcstas-comps/examples+ in the source
tree), organised by facility/category. Any of them can be loaded from
\texttt{mcgui}'s \textbf{New from Template} menu (known as ``Neutron site''
in older \MCS releases, see section~\ref{p:neutronsite}): selecting one
copies the instrument into your working directory, ready to edit, run, or
trace.

The list below is a curated tour, not an exhaustive catalogue -- new
example instruments are added continuously as facilities contribute
instrument descriptions. Use \verb+mcdoc+ (section~\ref{s:mcdoc-run}) or
browse the \verb+examples/+ directory directly for the full, current list.

%%%%%%%%%%%%%%%%%%%%%%%%%%%%%%%%%%%%%%%%%%%%%%%%%%%%%%%%%%%%%%%%%%%%%%%%%%%%%%%%
\section{A quick tour of instrument examples}
\label{s:quick-tour-instr}

\subsection{Templates}

The \texttt{Templates} category is the recommended starting point for new
instrument work: skeleton instruments to copy and adapt, covering the most
common instrument classes -- \texttt{template} and \texttt{template\_simple}
(bare instrument skeletons), \texttt{templateTAS} (triple-axis, including a
basic in-plane UB-matrix transformation for resolution estimation),
\texttt{templateTOF} and \texttt{templateDIFF} (time-of-flight and powder
diffraction), \texttt{templateSANS}/\texttt{templateSANS2}/\texttt{templateSANS\_MCPL}
and \texttt{templateSasView} (small-angle scattering, the latter using the
\texttt{SasView\_model} interface to SasView/sasmodels form factors),
\texttt{templateLaue} and \texttt{templateNMX}/\texttt{templateNMX\_TOF}
(single-crystal Laue diffraction), and \texttt{Demo\_shape\_primitives}
(illustrates the general-purpose shape/geometry components). Instrument
parameters are documented in each file's header (readable via
\verb+mcdoc+) and are designed to cover most instruments of the
corresponding class.

\subsection{Risoe}

Historical Ris\o\ National Laboratory instruments, including the TAS1
triple-axis spectrometer family described in detail in section
\ref{s:TAS1} below. TAS1 was decommissioned along with the DR3 reactor, but
its detailed McStas model remains one of the most thoroughly
validated example instruments in the distribution (see the model-vs-measurement
comparisons in section~\ref{data:TAS1}), and it constitutes the default
regression test series for the \MCS distribution.

\subsection{ISIS}

Instruments modelled on ISIS spallation-source beamlines, several using the
\texttt{ISIS\_moderator} source component (built from MCNP moderator
calculations): \texttt{ISIS\_CRISP} and \texttt{ISIS\_SANS2d} (reflectometry
and SANS), \texttt{ISIS\_HET}, \texttt{ISIS\_LET}, \texttt{ISIS\_MERLIN} and
\texttt{ISIS\_TOSCA\_preupgrade} (chopper spectrometers), \texttt{ISIS\_GEM}
and \texttt{ISIS\_OSIRIS} (powder diffraction and backscattering), and
\texttt{ISIS\_Prisma2}, a simple time-of-flight instrument with a RITA-style
analyser described in detail in section~\ref{s:PRISMA} below.

\subsection{ILL}

Cold and thermal guide models for a large number of ILL instruments,
several based on detailed technical drawings, e.g. \texttt{ILL\_H15},
\texttt{ILL\_H16}, \texttt{ILL\_H22} and \texttt{ILL\_H53} (guide systems)
feeding instrument-specific models such as \texttt{ILL\_H15\_IN6} (a
hybrid time-of-flight spectrometer with three monochromators and a Fermi
chopper; a usage example of the \texttt{GROUP} and \texttt{SPLIT} keywords,
see sections~\ref{s:instrdefs-extend-group} and~\ref{s:instrdefs-extend-enhance}
in the Component Manual), \texttt{ILL\_H15\_D11} and \texttt{ILL\_H512\_D22}
(SANS), \texttt{ILL\_H22\_D1A}/\texttt{ILL\_H22\_D1B} and \texttt{ILL\_D2B}
(powder diffraction), \texttt{ILL\_H53\_IN14} and \texttt{ILL\_H142\_IN12}
(triple-axis), \texttt{ILL\_IN5} and \texttt{ILL\_IN4} (time-of-flight),
and \texttt{ILL\_H22\_VIVALDI}/\texttt{ILL\_H143\_LADI} (Laue diffraction).

\subsection{ESS}

Design studies for the European Spallation Source, built around the
\texttt{ESS\_butterfly} moderator/source component, e.g.
\texttt{ESS\_IN5\_reprate} (a cold chopper multi-frame spectrometer for the
long-pulsed ESS timing structure) and \texttt{ESS\_BEER\_MCPL} (illustrating
MCPL-based event exchange between simulation stages, see the User Manual
section on MCPL/event files).

\subsection{Union\_demos, Union\_sample\_environments, Union\_validation}

Worked examples of the Union framework (contributed by Mads Bertelsen),
which separates sample/environment \emph{geometry} from \emph{physical
process} definitions and supports multiple scattering within arbitrarily
complex, nested sample-environment assemblies. \texttt{Union\_demos/Demonstration}
is a good first read: several powder samples in cans, a sample holder, and
a surrounding cryostat, illustrating the syntax and possibilities of the
Union components. See the Component Manual chapter on samples for the
underlying Union component reference.

\subsection{NCrystal}

Examples using the \texttt{NCrystal\_sample} component, interfacing the
external NCrystal library for physically detailed elastic and inelastic
scattering from single crystals, powders, and (to some extent) liquids,
based on material scattering kernels.

\subsection{Tests\_* categories}

A large set of \texttt{Tests\_sources}, \texttt{Tests\_optics},
\texttt{Tests\_samples}, \texttt{Tests\_monitors}, \texttt{Tests\_polarization},
\texttt{Tests\_union}, \texttt{Tests\_grammar}, \texttt{Tests\_RNG},
\texttt{Tests\_MCPL\_etc} and \texttt{Tests\_other} instruments exercise
individual components or language features in isolation. These are the
backbone of the \MCS regression test suite (run via \texttt{mctest}, see
section~\ref{s:mctest}), but are equally useful as minimal starting points
when learning how a specific component behaves. \texttt{Tests\_samples/Samples\_vanadium}
is described in detail in section~\ref{s:V-instr} below.

\subsection{Other facility categories}

Further contributed instrument models exist for BNL, DTU, FZ\_J\"ulich, HZB,
HighNESS, LLB, Necsa, PSI, SINE2020, SNS, TRIGA and TU~Delft, alongside a
\texttt{Mantid} category (instruments paired with Mantid IDF export via
\texttt{mcdisplay-mantid\_xml}, see section~\ref{s:mcdisplay}) and a
\texttt{Tools} category (utility instruments, e.g. for reading back
third-party event files via MCPL). Use \verb+mcdoc+ to search across all
categories by keyword, component used, or facility name.

%%%%%%%%%%%%%%%%%%%%%%%%%%%%%%%%%%%%%%%%%%%%%%%%%%%%%%%%%%%%%%%%%%%%%%%%%%%%%%%%
\section{A test instrument for the component V\_sample}
\label{s:V-instr}

This is one of many test instruments written with the
purpose of testing the individual components. We have picked
this instrument both to present an
example test instrument and because it despite its simplicity
has produced quite non-trivial results.

The instrument, \texttt{Tests\_samples/Samples\_vanadium}, consists of a narrow source,
a 60' collimator, a V-sample shaped as a hollow cylinder
with height 15 mm, inner diameter 16 mm, and outer diameter 24 mm
at a distance of 1 m from the source.
The sample is in turn surrounded by an unphysical $4\pi$-PSD
monitor with $50 \times 100$ pixels and a radius of $10^{6}$ m.
The set-up is shown in figure~\ref{f:V-instr}.

\begin{figure}
  \begin{center}
%FIXME    \psfrag{4pipsd}{$4\pi$ PSD}
    \includegraphics[width=0.9\textwidth]{figures/vanadium}
  \end{center}
\caption{A sketch of the test instrument for the component
V\_sample.}
\label{f:V-instr}
\end{figure}

\subsection{Scattering from the V-sample test instrument}
\label{s:vanadium-result}

In figure \ref{f:V-results}, we present the radial distribution of the
scattering from an evenly illuminated V-sample, as seen by a spherical PSD.  It
is interesting to note that the variation in the scattering intensity is as
large as 10\%. This is an effect of anisotropic attenuation of the beam in the
cylindrical sample.

\begin{figure}
  \begin{center}
    \includegraphics[width=0.6\textwidth]{figures/vanadium-surf-2}
  \end{center}
\caption{Scattering from a V-sample, measured by a spherical
  PSD. The sphere has been transformed onto a plane and the intensity is
  plotted as the third dimension. }
\label{f:V-results}
\end{figure}

%%%%%%%%%%%%%%%%%%%%%%%%%%%%%%%%%%%%%%%%%%%%%%%%%%%%%%%%%%%%%%%%%%%%%%%%%%%%%%%%
\section{The triple axis spectrometer TAS1}
\label{s:TAS1}
\index{TAS1}
\index{Triple-axis spectrometer!TAS1}

With this instrument definition, we have tried to create
a very detailed model of the conventional cold-source
triple-axis spectrometer TAS1 at the now closed neutron source DR3 of
Ris\o\ National Laboratory.
Except for the cold source itself, all components
used have quite realistic properties. Furthermore, the overall
geometry of the instrument has been adapted from
the detailed technical drawings of the real spectrometer.
The TAS 1 simulation was the first detailed work
performed with the \MCS package.
For further details see reference~\cite{tas1_report}.
The instrument family lives under \texttt{Risoe/} in the example library,
e.g. \texttt{Risoe/TAS1\_Vana} and \texttt{Risoe/TAS1\_Powder}, and is also
known by its historical name \texttt{linup-*} in older documentation.

At the spectrometer, the channel from the cold source
to the monochromator is asymmetric, since the first
part of the channel is shared with other instruments.
In the instrument definition, this is represented by
three slits.
For the cold source, we use a flat energy
distribution (component \textbf{Source\_flat})
focusing on the third slit.

The real monochromator consist of seven blades, vertically focusing on
the sample. The angle of curvature is constant so that the focusing is
perfect at 5.0 meV (20.0 meV for 2nd order reflections) for a 1$\times$1~cm$^2$
sample. This is modeled directly in the instrument definition using
seven \textbf{Monochromator} components. The mosaicity of the pyrolytic
graphite crystals is nominally 30' (FWHM) in both directions.  However, the
simulations indicated that the horisontal mosaicities of both
monochromator and analyser were more likely 45'. This was used for all
mosaicities in the final instrument definition.

The monochromator scattering angle, in effect determining the incoming
neutron energy, is for the real spectrometer fixed by four holes in the
shielding, corresponding to the energies 3.6, 5.0, 7.2, and 13.7~meV for
first order neutrons.  In the instrument definition, we have adapted the
angle corresponding to 5.0~meV in order to test the simulations against
measurements performed on the spectrometer.

The width of the exit channel from the monochromator may
be narrowed down from initially 40~mm
to 20~mm by an insert piece. In the simulations, we have chosen
the 20~mm option and modeled the channel with two slits to match
the experimental set-up.

In the test experiments, we used two standard samples:
An Al$_2$O$_3$ powder sample and a vanadium sample. The instrument
definitions use either of these samples of the correct
size. Both samples are chosen to focus on the opening aperture of
collimator 2 (the one between the sample and the analyser).
Two slits, one before and one after the sample,
are in the instrument definition set to the opening values which
were used in the experiments.

The analyser of the spectrometer is flat and made from
pyrolytic graphite. It is placed between an entry and
an exit channel, the latter leading to a single detector.
All this has been copied into the instrument definition.

On the spectrometer, Soller collimators may be inserted
at three positions: Between monochromator and sample,
between sample and analyser, and between analyser and detector.
In our instrument definition, we have used 30', 28', and 67' collimators
on these three positions, respectively.

An illustration of the TAS1 instrument
is shown in figure~\ref{f:TAS1}.
Test results and data from the real spectrometer are shown
in Appendix~\ref{data:TAS1}.

\begin{figure}
  \begin{center}
    \includegraphics[width=0.9\textwidth]{figures/tas1}
  \end{center}
\caption{A sketch of the TAS1 instrument.}
\label{f:TAS1}
\end{figure}

\subsection{Simulated and measured resolution of TAS1}
\label{data:TAS1}

In order to test the \MCS package on a qualitative level,
we have performed a very detailed comparison of a simulation with a
standard experiment from TAS1. The measurement series
constitutes a complete alignment of the spectrometer,
using the direct beam and scattering from V and Al$_2$O$_3$
samples at an incoming energy of 20.0~meV, using the second order
scattering from the monochromator.

In these simulations, we have tried to reproduce
every alignment scan with respect to position and width
of the peaks, whereas we have not tried to compare
absolute intensities. Below, we show a few comparisons
of the simulations and the measurements.

Figure \ref{f:2t_direct} shows a scan of
$2\theta_m$ on the collimated direct beam in two-axis mode.
A \hbox{1 mm} slit is placed on the sample position.
Both the measured width and non-Gaussian peak shape
are well reproduced by the \MCS simulations.

\begin{figure}[htp!]
  \begin{center}
    \includegraphics[width=0.6\textwidth]{figures/tas1-2T.pdf}
  \end{center}
\caption{TAS1: Scans of $2\theta_s$ in the direct beam with 1 mm slit on the
  sample position.
"$\times$": measurements, "o": simulations, scaled to the same intensity
Collimations: open-30'-open-open.}
\label{f:2t_direct}
\end{figure}

In contrast, a simulated $2\theta_a$ scan in triple-axis
mode on a V-sample showed a surprising offset from 0 degrees.
However, a simulation with a PSD
on the sample position showed that the beam center was 1.5~mm
off from the center of the sample, and this was important
since the beam was no wider than the sample itself.
A subsequent centering of the beam resulted in a nice
agreement between simulation and measurements.
For a comparison on a slightly different instrument
(analyser-detector collimator inserted),
see Figure~\ref{f:v_2ta_zero}.

\begin{figure}[htp!]
  \begin{center}
    \includegraphics[width=0.6\textwidth]{figures/vanadium-plot-2}
  \end{center}
\caption{TAS1: Corrected $2\theta_a$ scan on a V-sample.
Collimations: open-30'-28'-67'.
"$\times$": measurements, "o": simulations.}
\label{f:v_2ta_zero}
\end{figure}

The result of a $2\theta_s$ scan on an Al$_2$O$_3$
powder sample in two-axis mode is shown in Figure \ref{f:al2o3}.
Both for the scan in focusing mode (+ $-$ +)
and for the one in defocusing mode (+ + +) (not shown),
the agreement between simulation and experiment is excellent.

\begin{figure}[htp!]
  \begin{center}
    \includegraphics[width=0.6\textwidth]{figures/al2o3-focus}
  \end{center}
\caption{TAS1: $2\theta_s$ scans on Al$_2$O$_3$ in two-axis, focusing mode.
Collimations: open-30'-28'-67'.
"$\times$": measurements, "o": simulations.
A constant background is added to the simulated data.}
\label{f:al2o3}
\end{figure}

As a final result, we present a scan of the energy
transfer $E_a = \hbar \omega$ on a V-sample.
The data are shown in Figure \ref{f:v_ea}.

\begin{figure}[htp!]
  \begin{center}
    \includegraphics[width=0.6\textwidth]{figures/ea-scan}
  \end{center}
\caption{TAS1: Scans of the analyser energy on a V-sample.
Collimations: open-30'-28'-67'.
"$\times$": measurements, "o": simulations.}
\label{f:v_ea}
\end{figure}

For a modern triple-axis instrument that is not tied to a specific,
decommissioned facility, see the \texttt{Templates/templateTAS} example
instead, which additionally includes a basic in-plane UB-matrix
transformation useful for estimating resolution functions of an
arbitrary TAS configuration (see also \texttt{mcresplot}, section
\ref{s:mcresplot}).


%%%%%%%%%%%%%%%%%%%%%%%%%%%%%%%%%%%%%%%%%%%%%%%%%%%%%%%%%%%%%%%%%%%%%%%%%%%%%%%%
\section{The time-of-flight spectrometer PRISMA}
\label{s:PRISMA}
\index{PRISMA}
\index{Time-of-flight spectrometer!PRISMA}

In order to test the time-of-flight aspect of \MCS, we have
in collaboration with Mark Hagen, now at SNS, written a simple
simulation of a time-of-flight instrument loosely based on the ISIS
spectrometer PRISMA. The simulation was used to investigate the effect
of using a RITA-style analyser instead of the normal PRISMA backend.
The instrument is distributed as \texttt{ISIS/ISIS\_Prisma2} in the
example library.

We have used the simple time-of-flight source \textbf{Tof\_source}.
The neutrons pass through a
beam channel and scatter off from a vanadium sample, pass through
a collimator on to the analyser.
The RITA-style analyser consists of seven analyser crystals
that can be rotated independently around a vertical axis. After the
analysers we have placed a PSD and a time-of-flight detector.

To illustrate some of the things that can be done in a simulation as
opposed to a real-life experiment, this example instrument further
discriminates between
the scattering off each individual analyser crystal
when the neutron hits the detector. The
analyser component is modified so that a global variable
\verb+neu_color+ registers which
crystal scatters the neutron ray. The detector component
is then modified to construct seven different time-of-flight histograms,
one for each crystal (see the source code for the instrument
for details). One way to think of this is that
the analyser blades paint a color on each neutron which is then
observed in the detector. Since \MCS~3.x, the same effect could
alternatively be achieved more simply using a \texttt{USERVARS} field
(section~\ref{s:uservars}) rather than a global \texttt{DECLARE} variable,
which additionally makes the technique safe under SPLIT/GPU execution.
An illustration of the instrument is shown in figure~\ref{f:PRISMA}.
Test results are shown in Appendix~\ref{data:PRISMA}.

\begin{figure}[htp!]
  \begin{center}
    \includegraphics[width=0.9\textwidth]{figures/prisma2}
  \end{center}
\caption{A sketch of the PRISMA instrument.}
\label{f:PRISMA}
\end{figure}

\subsection{Simple spectra from the PRISMA instrument}
\label{data:PRISMA}

A plot from the detector in the PRISMA simulation is shown in Figure
\ref{f:PRISMAdata}. These results were obtained with each analyser blade
rotated one degree relative to the previous one. The separation of the
spectra of the different analyser blades is caused by different energy
of scattered neutrons and different flight path length from source to
detector.  We have not performed any quantitative analysis of the data.

\begin{figure}[htp!]
  \begin{center}
    \includegraphics[width=0.6\textwidth]{figures/prisma2-a}
  \end{center}
\caption{Test result from PRISMA instrument using ``colored
  neutrons''. Each graph shows the neutrons scattered from one analyser blade.}
\label{f:PRISMAdata}
\end{figure}
